% Report fragment for PA194AX11. Requires amsmath, amssymb, booktabs.
\section{PA194AX11: storing the gate and continuing through modus ponens}
\label{sec:ax11}

PA194AX11 changes the command interface, rather than only the naming of
states or the ordering of equivalent axiom formulas. The old AX command
left its guarded axiom on top and returned to the scheduler. The new AX
pushes that axiom onto the stack, then continues into the already available
modus-ponens routine. Successful modus ponens no longer pushes a dummy
zero. The resulting machine has 194 working states and 388 transitions,
with dense start zero and halt label 194 on the all-zero, two-way infinite
binary tape. There is no external initialization, additional tape, or
stationary-head move.

The mathematical target is
\[
 \mathrm{PA}\vdash
 \bigl(\neg\operatorname{Halt}_{\mathrm{blank}}(\mathrm{PA194AX11})
                \leftrightarrow\operatorname{Con}(\mathrm{PA})\bigr).
 \tag{AX11}
\]
The proof keeps the original coefficient-four separated-diagonal calculus.
The new work is the exact source splice, the four-parameter AX command,
a complete all-store proof compiler for its eleven-entry stack, and the
literal-machine binding. The logical calibration is the explicitly
identified inherited ordinary mathematical PA argument. No new
end-to-end Lean, Rocq, or iPRA endpoint is claimed.

\subsection{Two exact source splices}

In the principal PA194 source, the AX emitter ends in one
\texttt{breakTo 1} leaf. Replace that leaf by the existing stack-push
macro \texttt{node49}. The new macro has nineteen primitive sites and
falls through to the following full modus-ponens routine. Independently,
change \texttt{node163=(node6,node55)} to
\texttt{node163=(node6,node54)}, deleting the successful-MP dummy push.
Its native implementation changes transition $(164,1)$ from
$(1,R,98)$ to $(1,R,114)$.

The independently resolved source graph checks both splices completely.
Old ordinal 3553, the AX return leaf, becomes new push ordinals
3553 through 3571; its exit is the full MP field-reader entry 3572,
corresponding to old ordinal 3554. Old successful-MP push ordinals
3634 through 3652 are removed, and its retained copy block, old
3653 through 3659, moves to new 3652 through 3658. All 3,653 other
primitive instructions, including every resolved successor, agree under
the explicit relocation. No outside edge enters a splice interior.
The whole source has
\[
                 3673-1+19-19=3672
\]
primitive sites. These are complete finite graph equalities, checked by
two independently implemented source resolvers. The inserted nineteen-site
push is also exactly the old push block under its own entry/exit relocation;
all internal targets agree, and its normal exit is the MP reader entry.

The AX macro word still has 190 tokens. Its final three tokens are
now guard-copy, guard-comparison, and push, instead of guard-copy,
guard-comparison, and return. A free-constructor calculation checks all
six ordinary axiom heads, the restricted implication head, the diagonal
head, the entire inherited stack, and the computed gate. Both outcomes
of the guard are pushed: the guarded axiom when allowed, zero otherwise.
The old top is still saved. This adds an eleventh stack entry and leaves
the visible top zero.

\subsection{The extra operand is essential}

Continuing into the MP routine does \emph{not} automatically call
$\mathsf{MP}_0$. Clearing the AX parameter register does not clear the
continuation register. The first block of \texttt{node75} unpairs that
continuation, reads a new operand $z$, and replaces the continuation by
its tail. Thus AX has \emph{four} numerical parameters, written
$\mathsf A_{c,x,y,z}$.

After the added push, the visible command registers are $t=d=h=0$.
The ensuing no-push MP compares $I(0,z)$ with zero. If $z=0$ it succeeds
and returns zero; if $z>0$ it fails and returns zero. In both cases it
preserves the just-created stack, leaves $d=z$, and clears $h$.
Therefore the exact command effect is
\[
 \mathsf A_{c,x,y,z}(t,w,d,h)
       =(0,[q_0,q_1,\ldots,q_9,t;w],z,0).       \tag{A}
\]
Here $D=d+x$, $H=h+y$, and the entries are
\[
\begin{array}{c|l}
 j&q_j\\\hline
 0&G(H,D,c)\text{, or zero when the guard fails}\\
 1&\Gamma(c)\\
 2&h_6(D,c,H)\\
 3&h_5(D,H)\\
 4&h_4(D,c,H)\\
 5&h_3(D,c,H)\\
 6&h_2(D,c,H)\\
 7&h_1(D,H)\\
 8&\rho_8(D,c,H,d)\\
 9&U(d,h).
\end{array}
\]
The proof compiler below deliberately selects $z=0$, but soundness must
cover every $z$ because the machine enumerates all numerical candidates.
The extra field is therefore part of the command encoding and its
correctness proof, not an omitted precondition.

\subsection{The unchanged numerical calculus and its PA calibration}

For $p(a,b)=(a+b)(a+b+1)/2+b$, use
\[
 I(A,B)=4p(A,B),\quad M(a,b)=4p(a,b)+1,\quad
 U(a,B)=4p(B,a)+2,\quad E(a,b)=4p(a,b)+3.
\]
Write $A\to B$ for $I(A,B)$, $N(B)=I(B,1)$,
$A\wedge B=N(I(A,N(B)))$, and $\exists_a B=N(U(a,N(B)))$.
The six ordinary axiom families are
\begin{align*}
 h_1(D,H)&=(\exists_D U(D,H))\to H,\\
 h_2(D,A,H)&=U(D,U(A,H))\to U(A,U(D,H)),\\
 h_3(D,A,H)&=M(H,D)\to(E(H,A)\to M(A,D)),\\
 h_4(D,A,H)&=E(H,D)\to(M(A,H)\to M(A,D)),\\
 h_5(D,H)&=\exists_D E(H,D),\\
 h_6(D,A,H)&=U(D,A)\to(U(D,A\to H)\to U(D,H)).
\end{align*}
The restricted implication family is
\[
 \rho_8(D,A,H,u)=((H\to u)\to D)\to((D\to H)\to(A\to H)),
 \qquad u\leq D.
\]
There are also $\Gamma(B)=I(B,U(B,B))$ and the guarded gate below:
\begin{align*}
 V(a,b,c)&=M(a,b)\to U(c,M(a,b)),\\
 P(a,b,c)&=c\to M(c,b),\\
 S(a,b,c)&=U\bigl(b,(M(a,b)\to P(a,b,c))
                              \wedge(M(c,b)\to c)\bigr),\\
 Q(a,b,c)&=U(c,S(a,b,c)\to N(P(a,b,c))),\\
 G(a,b,c)&=N(V(a,b,c)\to Q(a,b,c)),\qquad c\ne a,b.
\end{align*}
The rules are modus ponens and theorem generalization $B\mapsto U(i,B)$.
There are no undischarged hypotheses, although labels may be free in a
formula. Call this numerical calculus $\mathcal C_4$. Zero is a theorem,
since $\rho_8(0,0,0,0)=0$, and one is the designated contradiction.
The code-zero identity $I(0,0)=0$ is a numerical alias, not a separate
syntactic constructor case.

The inherited PA195, diagonal, and separate-head reports establish
primitive-recursive refutation translations with PA-provable correctness:
\[
 \mathrm{PA}\vdash
 \left(\exists e\,\operatorname{Prf}_{\mathcal C_4}(e,1)
 \leftrightarrow
 \exists p\,\operatorname{Prf}_{\mathrm{PA}}
                         (p,\ulcorner0=1\urcorner)\right). \tag{C}
\]
For completeness, the retained membership-rewriting and equality-witness
interfaces, together with the diagonal replacement of equality vacuity,
recover the PA-refutation compiler. For the reverse direction, interpret
membership as the finite bit relation, equality as equality, and universal
codes as first-order universal quantification. Code one becomes
$\operatorname{Bit}(x_0,x_0)$, whose negation PA proves. The translation
of $I(0,0)$ uses the fixed equivalence of truth with truth-implies-truth.
Every occurring source label is smaller than its formula code, so the
self-indexed binder in $\Gamma(B)$ is fresh.

For each arithmetic formula $\theta$, finite binary comprehension has a
PA proof obtained by substituting $\theta$ into a fixed induction template.
The gate uses the matrix
$(\operatorname{Bit}(x_a,z)\vee z=x_b)\wedge\theta[z/x_b]$
with bound $x_a+x_b+1$ when $a\ne b$, and the matrix
$(z=x_b)\wedge\theta[z/x_b]$ with bound $x_b+1$ when $a=b$.
In the second case the bound membership antecedent is
$\operatorname{Bit}(z,z)$ and is false. The guard and code freshness
make the set witness safe. The proof supplier is primitive recursive in
the formula syntax; it is not a truth decision procedure or a reflection
principle. The displayed induction is full PA induction for arbitrary
formulas. Replacing this calculus by iPRA's smaller calculus would not
supply the same consistency calibration.

PA194AX11 leaves every formula family in (C) unchanged. Thus no new
recoding, new logical completeness theorem, or new treatment of the
dominated axiom is required for this construction. What remains is to
show that its changed command language enumerates precisely enough
$\mathcal C_4$ derivations and cannot manufacture a false theorem.

\subsection{Exact command words on arbitrary Cantor stacks}

The remaining commands, with $b=d+x$, are
\[
 \mathsf S(t,w,d,h)=
 \begin{cases}\mathsf{halt},&t+h=1,\\
 (h,\operatorname{tl}w,d,\operatorname{hd}w),&t+h\ne1,
 \end{cases}
\]
\[
 \mathsf{MP}_x(t,w,d,h)=
 \begin{cases}(b,w,b,0),&I(t,b)=h,\\
 (0,w,b,0),&I(t,b)\ne h.
 \end{cases}                                               \tag{M}
\]
The sum guard is retained exactly. In particular, it does not halt at
$(t,h)=(1,1)$. The stack pairing is total, with
$\operatorname{hd}(0)=\operatorname{tl}(0)=0$ and $[0;0]=0$.
We use these identities directly; stacks are not assumed to be free lists.

Let $A_0=\mathsf A_{0,0,0,0}$. All entries $q_0,\ldots,q_8$ have residue
zero modulo four, and $q_9=U(d,h)$ has residue two and is at least two.
Directly iterating the Shift equation gives, without halting,
\begin{align*}
 \mathsf A_{c,x,y,0};\mathsf S^{10}(t,w,d,h)
      &=(q_8,[t;w],0,q_9),\\
 \mathsf A_{c,x,y,0};\mathsf S^{11}(t,w,d,h)
      &=(q_9,w,0,t).
\end{align*}
Before the first ten moves, every tested sum is zero or a sum of
residue-zero entries. The next sum includes $q_9\geq2$, so it is also
safe. An additional Shift tests $q_9+t\geq2$, independently of $t$.

Define
\[
 \mathsf{Push}=A_0;\mathsf S^{10};\mathsf{MP}_0,\qquad
 \mathsf{Pop}=A_0;\mathsf S^{12},\qquad
 \mathsf{Discard}=\mathsf{MP}_0,\qquad
 \mathsf{Offset}_i=\mathsf{MP}_i.
\]
For all displayed natural values,
\begin{align*}
 \mathsf{Push}(t,w,d,h)&=(0,[t;w],0,0),\\
 \mathsf{Pop}(t,[a;w],d,h)&=(t,w,0,a),\\
 \mathsf{Discard}(t,w,0,h)&=(0,w,0,0),\\
 \mathsf{Offset}_i(0,w,0,0)&=(0,w,i,0).
\end{align*}
Push's comparison fails because implication and universal codes have
distinct residues. Discard's two MP branches coincide when the proposed
conclusion is zero. The Pop identity is unconditional, even at $t=1$;
its last test contains $q_9$ and cannot equal one. These equations also
cover an empty stack and every collapsed zero push.

For $0\leq j\leq8$, from any $(t,w,0,0)$ the selector
\[
 \mathsf{Discard};\mathsf A_{A,D,H,0};
            \mathsf S^{j+2};\mathsf{Pop}^{9-j}             \tag{Sel}
\]
returns $(q_j,w,0,0)$. After the initial Shifts, $q_j$ is the top,
$q_{j+1}$ the cache, and the remaining generated heads followed by the
saved zero form the stack above $w$. The indicated Pops remove precisely
that suffix while preserving the selected top. This includes the gate
$j=0$ and diagonal $j=1$ as well as all six ordinary heads. To select an
actual restricted leaf $\rho_8(D,A,H,u)$, with $u\leq D$, use
\[
 \mathsf{Discard};\mathsf{Offset}_u;
 \mathsf A_{A,D-u,H,0};\mathsf S^{10};\mathsf{Pop}.
\]
The offset is accumulated before AX and cleared by its explicitly zero
fourth parameter. Its emitted restricted head has the required actual
inequality $u\leq D$.

Generalization uses
\[
 \mathsf{Gen}_i=\mathsf{Push};\mathsf{Offset}_i;\mathsf S;
                \mathsf A_{i,0,0,0};\mathsf S^{11}.
\]
It sends $(a,w,0,0)$ to $(U(i,a),w,0,0)$ for every $a,w,i$, including
$a=1$ and $i=0$. The intervening Shift sees $(t,h)=(0,0)$; the possible
old contradiction is consumed as the cache parameter for the new
universal formula, not tested as a visible contradiction.

\subsection{A proof compiler that works from every carried store}

On the input $(0,[a,I(a,b),0;w],0,0)$ the suffix
\[
                 \mathsf S;\mathsf S;\mathsf{MP}_b;\mathsf{Pop}
                                                               \tag{MPs}
\]
halts at the second Shift if $a=1$, and otherwise returns $(b,w,0,0)$.
This includes $b=1$: Pop is safe, and returns that contradiction for a
later final Shift. The zero below the two premises is essential because
the MP routine itself no longer pushes it. When $a=b=w=0$, all relevant
pairs and the implication can collapse to zero; the same calculation
uses their total projections and still returns $(0,0,0,0)$.

Compile each numerical-calculus proof $e$ into a finite word $W_e$.
Use (Sel) and the restricted selector for axiom leaves, and append
$\mathsf{Gen}_i$ for a generalization node. At an MP node with major
proof $e_1$ of $I(a,b)$ and minor proof $e_0$ of $a$, use
\[
 W_e=\mathsf{Discard};\mathsf{Push};W_{e_1};\mathsf{Push};
       W_{e_0};\mathsf{Push};\mathsf S;\mathsf S;
                              \mathsf{MP}_b;\mathsf{Pop}.
\]
Induction on the finite proof tree proves the precise invariant:
from every $(t,w,0,0)$, $W_e$ either halts or returns
$(\operatorname{concl}(e),w,0,0)$. The first Discard and Push in the
MP case reserve the zero before the subproofs; each subproof preserves
the entire protected old stack. Earlier halts satisfy the invariant
and require no special consistency assumption.

If $e$ proves one, the word
\[
                             A_0;W_e;\mathsf S             \tag{H}
\]
halts from \emph{every} store $(t,w,d,h)$. Its initial AX sets both
caches to zero while preserving some finite stack, to which the induction
hypothesis applies. If no earlier halt occurs, its last configuration is
$(1,w',0,0)$ and the sum guard halts. The final command is a plain Shift;
using the new Pop as the terminal test would be wrong, since Pop preserves
a visible one without halting.

\subsection{Four-field scheduling and soundness}

The original scheduler's fetch loop is retained. Its source now decodes
one additional nested operand on an AX command. Define $e(\epsilon)=0$
and $e(c::w)=1+b_c(e(w))$, where
\[
\begin{aligned}
 b_{\mathsf S}(z)&=p(0,z),\\
 b_{\mathsf{MP}_x}(z)&=p(1,p(x,z)),\\
 b_{\mathsf A_{c,x,y,u}}(z)&=p(c+2,p(y,p(x,p(u,z)))).
\end{aligned}                                               \tag{Enc}
\]
The four-parameter order is $c$ in the command tag, followed by fields
$y,x,u$. Every displayed inverse identity follows from total Cantor
projections, including all zero aliases. Continuing to use the old
three-parameter encoding would cause the new MP reader to consume part
of the next command. The independent command checker explicitly detects
that incorrect encoding.

For a positive continuation $1+b_c(z)$, its next continuation is a
component of $b_c(z)$ and is therefore at most $b_c(z)$, hence strictly
smaller than the positive continuation. The extra field does not change
this argument. Each arithmetic macro terminates with a primitive-recursive
bound on its input registers, so every nonhalting candidate finishes in
finite time. The fresh candidates begin at two and increase successively.

If $z=e(A_0;W_e;\mathsf S)$ for a refutation word (H), then prefixing
$\mathsf A_{k,0,0,0}$ gives fresh candidate codes
\[
                 p(k+2,p(0,p(0,p(0,z)))) .                  \tag{Cof}
\]
They form an unbounded family, all at least three. The prefix can alter
the carried store, but the suffix halts from every store. PA proves by
bounded induction on the finitely many earlier candidates that one of
these halting candidates is reached unless the machine halts earlier.
Thus a calculus refutation gives a scheduler halt. No reset of the stack
between trials is assumed.

For soundness annotate the top, cache, and every canonical Cantor-stack
entry with a $\mathcal C_4$ proof. Zero has its fixed proof. Every new AX
entry in (A), including its pushed gate and saved old top, is a theorem;
the post-AX MP puts zero in both visible positions for all $u$, and the
arbitrary offset $u$ carries no proof obligation. Successful MP combines
annotated premises, failure returns zero, and Shift moves annotations.
If $t+h=1$, nonnegativity implies that one visible entry is one, yielding
a calculus refutation. This argument does not incorrectly infer a halt
from $t=1$ when $h=1$. Positive Cantor tails decrease strictly, so the
finite canonical stack and its annotations have primitive-recursive
codes, including the virtual zero entries.

All these constructions and inductions are available in PA. Combining
the sound-store invariant, compiler, scheduling argument, and (C) yields
\[
 \mathrm{PA}\vdash
  H_{\mathrm{AX11\ source}}
   \leftrightarrow
    \exists p\,\operatorname{Prf}_{\mathrm{PA}}(p,\ulcorner0=1\urcorner).
\]

\subsection{From source to the exact blank-tape machine}

The new source has 3,672 sites and 4,779 source-outcome contracts. The
construction checker verifies every outcome for arbitrary registers and
both decrement branches, using 1,479,769 symbolic operations with maximum
case 497. The independent physical checker uses the literal transition
table, derives its four generic register-list lemmas from those rows, and
checks all source addresses using two independent successor resolvers.
The full finite source relocation and eleven-entry algebra are checked
separately; primitive simulation alone would not identify the new
four-parameter command semantics.

Exactly 58,154 literal steps from blank reach the clean source entry 25
at native state 84 and head $-2$, with logical registers
$(2,0,0,0,0,0,0,0,0)$ and 26 unit spare registers. The startup test compares
the complete finite-support tape and its blank exterior. The machine is
not supplied with that frame as nonblank input.

Each physical primitive simulation has positive duration and does not
halt prematurely. A source halt therefore yields a literal blank-tape
halt by finite concatenation after startup. Conversely, for a proposed
first literal halt at time $T$, if the source continued without halting,
concatenating $T+1$ positive source simulations would produce the same
deterministic literal run extending beyond $T$ without a halt. PA proves
this finite contradiction; no unbounded search for a next source boundary
or inverse decoding of arbitrary intermediate tapes is needed.

The exact local table checks, unary-run inductions, register-list
inductions, and finite concatenations are formalizable in PA. Composing
this machine/source transfer with the preceding calibration proves
(AX11). The executable certificates substantiate the new finite
bindings and universal symbolic contracts. They are not presented as
a newly completed proof-assistant formalization of the PA biconditional.
Actual nonhalting follows under consistency of PA; no finite run test
establishes that consistency assumption.

\subsection{A concrete difference in command behavior}

The changed interface is observable even on the all-zero command store
with zero AX parameters. The old guard fails, so the old AX returns
$(0,W_{10},0,0)$, where $W_{10}$ is its ten-entry stack. This stack is
positive: among its entries is $U(0,0)=2$. The new padded AX returns
$(0,p(0,W_{10}),0,0)$. For every $w>0$, Cantor pairing satisfies
$p(0,w)>w$, so these outputs are different even though the visible top
and caches agree. Thus the extra push is not erased by the alias
$p(0,0)=0$. The new command also consumes a fourth field and can set the
offset to any supplied value. These are changes of recurring source
semantics. The separate action-graph census rules out disguising the
new literal table as an earlier table by state renaming or cosmetic
symmetries; neither statement claims a lower bound on possible state counts.
